\begin{frame}{ML1 총정리: 한 학기를 14장으로}
  선형/로지스틱회귀의 경사하강법에서 시작해, 생성적 분류(나이브
  베이즈/GDA), 거리·마진 기반(kNN, SVM), 트리 앙상블(GBDT),
  신경망(역전파, CNN), 시퀀스·어텐션·LLM, 비지도학습(PCA,
  임베딩), 잠재변수 생성모델(EM/GMM, VAE, GAN, Diffusion)까지.
  \vskip0.2em
  \scriptsize
  \begin{tabular}{@{}p{2.0cm}p{5.5cm}p{7.4cm}@{}}
    \toprule
    장 & 배운 것 & 핵심 질문 \\
    \midrule
    Ch02 & 선형·로지스틱회귀 & 연속값과 확률을 어떻게 예측하는가 \\
    Ch03 & 나이브베이즈/GDA & 데이터가 어떻게 생성됐는지부터
      거꾸로 분류할 수 있는가 \\
    Ch04 & kNN/k-means/GMM·EM & ``가깝다''만으로 예측·군집화,
      그 한계는 \\
    Ch05 & SVM & 확률이 아니라 \textbf{마진}으로 경계를 정할 수
      있는가 \\
    Ch06 & 정규화/모델선택 & \textbf{과적합}을 손실함수로 어떻게
      통제하는가 \\
    Ch07 & 트리/RF/GBDT/SHAP & 잘 나누는 질문을 어떻게 고르고,
      어떻게 \textbf{설명}하는가 \\
    Ch09 & 신경망/역전파/학습기법 & 층을 어떻게 학습시키고, 깊은
      망은 왜 어려운가 \\
    Ch10 & CNN 기초와 응용 & 이미지의 \textbf{국소 구조}를 어떻게
      효율적으로 다루고 재사용하는가 \\
    Ch11 & 시퀀스 모델 & \textbf{순서}가 있는 데이터를 어떻게
      기억·처리하는가 \\
    Ch12 & 어텐션/트랜스포머 & 순차 처리 \textbf{없이} 모든 단어
      의 관계를 한 번에 계산 \\
    Ch13 & 그래프 표현학습 & Random Walk/Node2Vec/PageRank —
      그래프를 어떻게 \textbf{벡터}로 바꾸는가 \\
    Ch14 & PCA/word2vec/임베딩 & \textbf{정답 없이} 구조를 찾고,
      벡터를 실전에 어떻게 쓰는가 \\
    Ch15 & VAE/GAN/Diffusion & \textbf{잠재변수}로부터 어떻게 새
      데이터를 만드는가 \\
    Ch17(부록) & LLM: 사전학습/프롬프팅/정렬 & 다음 토큰 예측이
      폭넓은 능력이 되고, 사람 뜻에 맞게 다듬기 \\
    \bottomrule
  \end{tabular}
\end{frame}
